\documentclass[conference]{IEEEtran}
\IEEEoverridecommandlockouts
\usepackage{amsmath,amssymb}
\usepackage{booktabs}
\usepackage{graphicx}
\usepackage{siunitx}
\usepackage{url}
\usepackage[hidelinks]{hyperref}
\usepackage{cite}
\usepackage{xcolor}

\title{When Does Future Mobility Information Have Packet-Level Value in Vehicular Optical Scheduling?}
\author{\IEEEauthorblockN{Panagiota Grosdouli}\IEEEauthorblockA{Department of Electrical and Computer Engineering\\School of Engineering\\Democritus University of Thrace\\Xanthi, Greece\\panagros1@ee.duth.gr}}

\begin{document}
\maketitle

\begin{abstract}
Mobility forecasts can anticipate the evolution of a vehicular optical link, yet their usefulness ultimately depends on whether the scheduler can convert that information into better packet service. We examine this question in a model-based PC-FMCW/DPSK vehicular optical system with causal motion prediction, time-varying link estimates, explicit queues, and hard packet deadlines. A guarded predictive scheduler is compared with Reactive Greedy using four pre-specified operating regimes and 20 paired holdout seeds per regime. The guard is selected independently on development seeds before holdout evaluation. The selected ratio, 1.0, prevents forecast-driven reordering when it would reduce the best available immediate service score. Relative to Reactive Greedy, mean goodput differences are 0.00000, +0.00035, -0.00050, and +0.00080 Mbps for the 0.1 s deadline, 0.5 s deadline, offered-load 1.1, and +3 dB SNR regimes, respectively. None establishes a practically meaningful positive or negative effect under the pre-specified analysis, while exact paired ties occur in 80--100\% of episodes. Development ablations show that relaxing the guard increases predictive reordering but introduces losses under the shortest deadline. The results identify an actionability bottleneck: future mobility information yields packet-level value only when it changes a feasible service decision whose benefit exceeds the cost of foregoing immediate service.
\end{abstract}

\begin{IEEEkeywords}
vehicular optical communications, predictive scheduling, trajectory prediction, PC-FMCW, DPSK, visible light communication, packet deadlines
\end{IEEEkeywords}

\section{Introduction}
The quality of a vehicular optical link can change rapidly as relative position, orientation, field of view, and line-of-sight geometry evolve \cite{memedi2021vehicular}. Phase-coded FMCW laser-headlamp research also motivates optical platforms in which sensing, illumination, and communication coexist \cite{liu2026laserheadlamp,kumbul2023smoothed,kumbul2023mimo}. These characteristics make vehicle motion relevant to scheduling: information about where a receiver is heading can indicate that a currently usable link is about to deteriorate or disappear.

Predicting that evolution, however, does not by itself improve communication performance. A scheduler must decide whether acting on a forecast is preferable to serving the best opportunity available now. This distinction becomes particularly important when packets have finite deadlines. Advancing service for a receiver whose link is expected to deteriorate may preserve a future delivery opportunity, but the same decision can delay a packet in another queue, sacrifice a stronger current transmission, or consume the limited slack of a more urgent packet.

The individual ingredients of this problem are well established. Trajectory prediction has been studied for dynamic narrow-field-of-view VLC \cite{jiang2020trajectory}; resource allocation in VLC-based vehicular networks has considered connectivity, load balancing, and network lifetime \cite{msongaleli2020resource}; and predicted vehicle motion has been incorporated into LTE-V2V resource selection and vehicular computation and resource allocation \cite{hajrasouliha2021trajectory,guo2022trajectoryoffload,lu2024predictive}. Low-latency and reliable vehicular resource allocation has also been investigated extensively \cite{guo2019latency}. What remains less clear is how much of the information contained in a mobility forecast survives the sequence of decisions between prediction and packet delivery.

We address that question at the packet level. At each scheduling instant, a causal motion history is mapped to predicted geometry and optical-link evolution. The scheduler combines those estimates with the current queue and deadline state, while realized future positions remain unavailable to the policy. The resulting decisions are evaluated against Reactive Greedy using paired episodes and a service guard selected on development data before the holdout is examined.

The study makes three contributions. First, it provides a cross-layer formulation that connects causal motion prediction to link evolution, scheduling decisions, and deadline-constrained packet outcomes. Second, it evaluates the incremental goodput of a guarded predictive scheduler under a paired development/holdout protocol rather than inferring communication value from forecast accuracy alone. Third, it uses guard ablations and episode-level ties to examine why additional future information may have little downstream effect. The resulting picture is one of limited actionability: protecting immediate service prevents harmful predictive reordering, but can simultaneously leave the scheduler with few consequential opportunities to exploit the forecast.

\section{Related Work}
\subsection{Vehicular optical communication}
Vehicular VLC couples communication performance closely to geometry, mobility, line of sight, and the directional characteristics of vehicle lighting \cite{memedi2021vehicular}. Msongaleli and Kucuk formulate resource utilisation in VLC-based VANETs in terms of network lifetime, connectivity, and load balancing, using an ILP formulation and a heuristic for denser networks \cite{msongaleli2020resource}. These studies establish the mobility-sensitive nature of vehicular optical resource allocation. The present work considers a different layer of the problem: whether knowledge of future motion changes packet service once queues and deadlines are explicitly represented.

\subsection{Trajectory information in communication decisions}
Jiang \emph{et al.} use trajectory prediction to track a target light source in a dynamic VLC system with a narrow field of view \cite{jiang2020trajectory}. In RF vehicular networks, Hajrasouliha and Shahgholi Ghahfarokhi use predicted vehicle locations for geo-based LTE-V2V resource selection and report improvements in packet-reception and blocking-related metrics in their evaluated scenarios \cite{hajrasouliha2021trajectory}. Prediction has also been coupled with task offloading and resource allocation in vehicular edge systems \cite{guo2022trajectoryoffload,lu2024predictive}. Together, this literature establishes that mobility forecasts can inform communication and computation decisions. Our concern is the subsequent question of whether such information remains useful after it is filtered through a deadline-constrained packet scheduler.

\subsection{From forecast quality to packet utility}
Trajectory forecasting is itself a mature estimation problem; interacting-multiple-model approaches, for example, combine alternative motion hypotheses to improve vehicle-state prediction \cite{gao2023imm}. A communication scheduler nevertheless operates one stage further downstream. In the system considered here, information passes through the sequence
\[
\text{motion history}\rightarrow\text{forecast}\rightarrow\text{predicted link state}\rightarrow\text{scheduling action}\rightarrow\text{packet outcome}.
\]
An accurate estimate at the beginning of this chain need not alter its final outcome. The scheduler may already prefer the same receiver from current information, or the queue and deadline state may make a forecast-driven alternative too costly. We therefore evaluate mobility prediction through the packet decisions it can change, rather than treating prediction error as a surrogate for communication utility.

\section{System Model}
At decision slot $t$, the scheduler has access to vehicle motion histories through $t$, packet queues, packet deadlines, and previous service outcomes. A causal predictor generates target positions over a horizon $H$. Ego-relative range and bearing are then mapped through an analytical optical-link model to predicted SNR, BER/PER, goodput, outage, and link-lifetime quantities. Realized packet outcomes use the geometry observed at the actual transmission slot. Future realized positions are not available to deployable policies.

The physical-layer setting follows the PC-FMCW/DPSK premise and reference parameters motivated by phase-coded FMCW waveform and transceiver studies \cite{kumbul2023smoothed,kumbul2023mimo}. The laser-headlamp architecture of \cite{liu2026laserheadlamp} further motivates the integration of optical sensing and communication in a vehicle-lighting platform. The end-to-end geometry-to-link mapping used in this study is analytical; absolute received power is not calibrated against measured vehicle-to-vehicle optical power.

Packets arrive in per-receiver queues and may be retransmitted following unsuccessful service attempts. For a physical deadline of $N$ slots, a packet has exactly $N$ service opportunities, and its final legal service index is $arrival+N-1$. Deadlines that are not exactly representable by the simulator slot duration are excluded rather than rounded. Queue accounting distinguishes successful delivery, expiration, retransmission, drops, and right-censored packets.

\section{Predictive Scheduling}
For receiver $i$, the predictive scheduler evaluates
\begin{equation}
 i^{\star}=\arg\max_i U_i\!\left(\mathcal{H}_t,Q_t,D_t,\widehat{\mathcal{L}}_{i,t+1:t+H}\right),
\end{equation}
where $\mathcal{H}_t$ is the causal motion history, $Q_t$ and $D_t$ describe queue and deadline state, and $\widehat{\mathcal{L}}_{i,t+1:t+H}$ denotes predicted link evolution. A receiver whose usable optical link is expected to disappear can therefore receive additional priority before the loss of connectivity occurs.

Forecast-driven urgency competes directly with the value of immediate service. To control that trade-off, a current-service guard allows the predictive scheduler to reorder receivers only when the candidate retains a prescribed fraction of the immediate service score obtained by the reactive choice. Guard ratios 0.8, 0.9, 0.95, and 1.0 are considered during development. The selected ratio is 1.0. Under this setting, a forecast-driven choice cannot replace the reactive choice with a receiver having a lower immediate service score according to the guard definition.

Random, Round Robin, Reactive Greedy, and Proportional Fair provide reactive scheduling references. Causal motion references include Last Position, constant velocity, constant acceleration, Kalman-CV, and IMM-style forecasting where applicable. Diagnostics based on future positions are confined to evaluator-side information references; they are neither deployable policies nor globally optimal offline schedulers.

\section{Experimental and Statistical Protocol}
The simulator uses a slot duration of 0.1 s, and packet deadlines are required to correspond to an integer number of slots. During validation, the deadline implementation was found to include one additional service opportunity beyond the intended $N$ slots. The implementation was corrected before the experiments reported here. A previously considered 0.05 s deadline is not representable at the chosen slot duration and is therefore excluded from the analysis.

Scheduler development uses seeds 20270101--20270110. Guard selection is completed on these seeds before evaluation on the disjoint holdout seeds 20270201--20270220. Four primary operating regimes are specified: a 0.1 s deadline, a 0.5 s deadline, offered load 1.1, and reference SNR +3 dB. The primary endpoint is the paired candidate-minus-Reactive-Greedy goodput difference. Statistical inference is performed at the paired scenario/episode level rather than over individual packets or time samples.

For each regime, the mean paired difference is estimated with 100,000 paired percentile-bootstrap replicates and a 95\% confidence interval. A two-sided paired Wilcoxon signed-rank test is used together with Holm correction across the four primary comparisons. We also report paired Cohen $d_z$ and the fractions of paired wins, losses, and ties. The pre-specified practical margin is $m=0.001$ Mbps. A positive effect requires the lower confidence-interval bound to exceed $+m$, whereas a negative effect requires the upper bound to fall below $-m$. Intervals that do not meet either condition are treated as inconclusive with respect to a practically meaningful effect. A non-significant hypothesis test is not interpreted as evidence of equality.

\section{Results}
\subsection{Incremental goodput on the holdout}
Development selection yields a guard ratio of 1.0. Table~\ref{tab:primary} reports the paired holdout comparison with Reactive Greedy. None of the four regimes satisfies the pre-specified criterion for a practically meaningful positive or negative effect.

\begin{table*}[t]
\centering
\caption{Primary paired comparisons against Reactive Greedy (20 paired holdout seeds per regime).}
\label{tab:primary}
\begin{tabular}{lrrrrl}
\toprule
Regime & $\Delta$ goodput (Mbps) & 95\% bootstrap CI & Holm $p$ & Cohen $d_z$ & Interpretation \\
\midrule
Deadline 0.1~s & +0.00000 & [0.00000, 0.00000] & 1.000000 & 0.000 & No observed difference \\
Deadline 0.5~s & +0.00035 & [0.00000, 0.00105] & 1.000000 & 0.224 & No established effect \\
Offered load 1.1 & -0.00050 & [-0.00455, 0.00285] & 1.000000 & -0.060 & No established effect \\
Reference SNR +3~dB & +0.00080 & [-0.00080, 0.00255] & 1.000000 & 0.208 & No established effect \\
\bottomrule
\end{tabular}
\end{table*}

At the 0.1 s deadline, all 20 paired episodes produce exactly the same goodput. At 0.5 s, 95\% of pairs tie and 5\% favor the predictive scheduler. For offered load 1.1, the corresponding fractions are 85\% ties, 10\% predictive wins, and 5\% Reactive-Greedy wins. At +3 dB, 80\% tie, 15\% favor the predictive scheduler, and 5\% favor Reactive Greedy. Thus, in most holdout episodes, adding the mobility forecast does not change realized goodput.

The all-zero contrast at the 0.1 s deadline makes the Wilcoxon statistic degenerate; its raw test is undefined and the Holm-adjusted value is conservatively assigned $p=1$. Raw Wilcoxon $p$ values for the remaining regimes are 0.3173, 1.0, and 0.2733, and all Holm-adjusted values equal 1.0. These comparisons do not establish a goodput advantage or loss. Except for the exact 0.1 s ties observed in this sample, they also do not establish equality between the schedulers.

\subsection{The service-order trade-off}
The development ablations clarify why the selected scheduler behaves so similarly to Reactive Greedy. With a guard of 0.8, mean development goodput increases in the 0.5 s and +3 dB regimes, but decreases under the 0.1 s deadline. That short-deadline loss violates the pre-specified worst-regime noninferiority requirement. Guards 0.9 and 0.95 also fail this requirement, leaving 1.0 as the only eligible setting before holdout evaluation.

Relaxing the guard gives predicted link urgency more influence over service order. This creates opportunities to preserve links that are expected to deteriorate, but it also permits the scheduler to give up a better immediate transmission. The 1.0 guard removes that immediate-service sacrifice and, at the same time, sharply restricts the set of decisions that prediction can change. The 80--100\% episode-level tie rates are consistent with this explanation: the scheduler frequently receives additional information without obtaining a consequential alternative action.

\subsection{Interpreting the value of future information}
The holdout evidence does not establish an incremental goodput advantage for the selected predictive scheduler. It also does not support the broader conclusion that mobility prediction has no communication value. The development results show that predictive reordering can affect packet outcomes, but its direction varies with operating conditions. Under the cross-regime robustness requirement used here, avoiding the observed short-deadline loss leaves a policy whose holdout goodput is usually indistinguishable at the episode level from that of Reactive Greedy.

This behavior locates the bottleneck between information and action. Future geometry can alter a link forecast, yet that change matters only if it modifies a feasible scheduling decision. The modified decision must then improve packet delivery enough to offset the service opportunity relinquished elsewhere. When either step fails, predictive information does not propagate to packet-level utility.

\section{Discussion}
\subsection{Forecast accuracy and communication value}
Trajectory-prediction studies typically assess geometric forecast quality, whereas communication and resource-allocation studies evaluate downstream system objectives \cite{gao2023imm,hajrasouliha2021trajectory,guo2022trajectoryoffload,lu2024predictive}. The present experiment illustrates why the two should be evaluated separately. Forecast accuracy determines the quality of information available to the scheduler; packet utility depends additionally on whether the scheduler can act on that information and whether the resulting action improves service.

The high frequency of paired ties is particularly informative in this respect. The predictive policy is not merely being evaluated on a small average improvement that fails a significance test. In many episodes it reaches exactly the same realized goodput as the reactive policy. For the selected guard, the decision rule leaves limited scope for future-link information to alter service without sacrificing immediate score. This observation motivates direct measurement of prediction-induced action changes in future scheduler designs.

\subsection{Link urgency versus packet urgency}
A predicted loss of connectivity creates link urgency, whereas a packet approaching expiration creates packet urgency. These two pressures need not identify the same receiver. Serving a link early because it may soon disappear can be beneficial when the current queue has sufficient slack elsewhere; the same action can be harmful when another receiver offers a stronger immediate transmission or contains a packet close to expiration. The guard-selection pattern is consistent with such a conflict, although the current experiment does not isolate a universal quantitative relationship between link urgency and packet urgency.

\subsection{Implications for predictive scheduler design}
The development and holdout results point to a design requirement that is separate from improving the motion predictor itself. A useful predictive scheduler must identify states in which future information changes a feasible decision and the expected benefit of that change justifies its immediate opportunity cost. A fixed conservative guard protects current service, but can make the predictive policy effectively reactive. A less restrictive rule creates more opportunities to exploit forecasts, while also exposing the system to harmful reordering. This trade-off suggests that future designs should make actionability state dependent, for example by accounting jointly for forecast confidence, expected link-opportunity loss, queue state, and remaining deadline slack.

\section{Limitations}
The evidence in this study is simulation based. The optical channel relies on analytical geometry and reference-SNR assumptions rather than measured end-to-end vehicular optical calibration. No road deployment, measured optical channel, or hardware-in-the-loop experiment is included. The cited PC-FMCW/DPSK literature motivates the physical-layer setting and parameters but does not validate the complete geometry-to-packet simulator used here.

The evaluation covers a finite scheduler family, four primary operating regimes, and 20 paired holdout seeds per regime. The selected 1.0 guard should therefore not be interpreted as globally optimal. The study does not provide a dense operating-region map, a calibrated sensitivity analysis over forecast error, or a globally optimal oracle upper bound. Moreover, the proposed actionability mechanism is inferred primarily from guard selection, development ablations, and episode-level tie structure; the present experiment does not directly decompose every prediction-induced scheduling change into beneficial and harmful actions.

Broader evaluation with measured optical channels, additional mobility and traffic conditions, learned trajectory models, larger driving datasets, and direct instrumentation of scheduler action changes would strengthen external validity and allow the actionability hypothesis to be tested more directly.

\section{Conclusion}
Future mobility information can influence packet delivery only through decisions that the scheduler is able and willing to change. In the vehicular optical system studied here, cross-regime development selection favors a 1.0 current-service guard, which prevents forecast-driven reordering from reducing the best immediate service score. On a disjoint 20-seed holdout, the resulting policy establishes no practically meaningful goodput advantage or loss relative to Reactive Greedy in any of the four primary regimes, and 80--100\% of paired episodes are exact ties. Development ablations show the underlying trade-off: relaxing the guard creates more forecast-driven opportunities, but also introduces a loss under the shortest deadline. These results place the limiting factor at the interface between prediction and scheduling. Future mobility information becomes packet-level value when it enables a consequential feasible action whose realized benefit is greater than the immediate queue and deadline cost of taking it.

\section*{Reproducibility}
The experiments use development seeds 20270101--20270110 and disjoint holdout seeds 20270201--20270220. The repository contains the simulator, experimental configuration, and analysis required to reproduce the reported paired comparisons. The results in this paper use the deadline semantics described in Section V; experiments generated with the earlier deadline implementation are excluded from the analysis.

\bibliographystyle{IEEEtran}
\bibliography{references}
\end{document}
